\chapter{Railway Design Rules}

\section*{Motivation}

\begin{remark}
Railway alignment design is not solely a geometric or mathematical task.
It is governed by engineering rules, safety requirements, and operational
constraints.
\end{remark}

\begin{remark}
These rules originate from:
\begin{itemize}
\item national regulations,
\item international standards (e.g.\ UIC),
\item operator-specific guidelines (e.g.\ Deutsche Bahn),
\item empirical engineering practice.
\end{itemize}
\end{remark}

\begin{remark}
Within the AIM framework, such rules are not treated as external
constraints only, but are integrated into the modelling and optimization
process.
\end{remark}

---

\section{Categories of Rules}

\subsection{Geometric Rules}

\begin{itemize}
\item minimum radius,
\item continuity of alignment,
\item transition length requirements,
\item curvature monotonicity in transitions.
\end{itemize}

\subsection{Kinematic and Dynamic Rules}

\begin{itemize}
\item limits on lateral acceleration,
\item limits on jerk,
\item compatibility of curvature and cant,
\item speed-dependent requirements.
\end{itemize}

\subsection{Operational Rules}

\begin{itemize}
\item allowable speeds,
\item safety margins,
\item maintenance considerations,
\item interoperability constraints.
\end{itemize}

\subsection{Topological Rules}

\begin{itemize}
\item constraints at switches and crossings,
\item parallel track spacing,
\item network connectivity conditions.
\end{itemize}

---

\section{Geometric Constraints}

\subsection{Minimum Radius}

\begin{definition}[Minimum radius constraint]
The curvature must satisfy
\[
|\kappa(s)| \leq \kappa_{\max},
\qquad
\kappa_{\max} = \frac{1}{R_{\min}}.
\]
\end{definition}

\begin{remark}
The value of \(R_{\min}\) depends on:
\begin{itemize}
\item design speed,
\item vehicle characteristics,
\item regulatory standards.
\end{itemize}
\end{remark}

---

\subsection{Continuity Conditions}

\begin{remark}
Railway alignments must satisfy at least:
\begin{itemize}
\item $C^0$ continuity (position),
\item $C^1$ continuity (direction),
\item typically $C^2$ continuity (curvature).
\end{itemize}
\end{remark}

\begin{remark}
Higher-order smoothness may be required for high-speed lines.
\end{remark}

---

\subsection{Transition Length}

\begin{definition}[Minimum transition length]
A transition between curvatures must satisfy
\[
L \geq L_{\min}(v,\kappa_1,\kappa_2).
\]
\end{definition}

\begin{remark}
Typical formulations relate transition length to:
\begin{itemize}
\item allowable rate of change of lateral acceleration,
\item cant development limits.
\end{itemize}
\end{remark}

---

\section{Dynamic Constraints}

\subsection{Lateral Acceleration}

\begin{definition}[Lateral acceleration]
The lateral acceleration is given by
\[
a_y(s) = v^2 \kappa(s).
\]
\end{definition}

\begin{constraint}
\[
|a_y(s)| \leq a_{\max}.
\]
\end{constraint}

---

\subsection{Cant and Cant Deficiency}

\begin{definition}[Cant angle]
Let $\phi(s)$ denote the cant angle.
\end{definition}

\begin{definition}[Effective acceleration]
The effective lateral acceleration is
\[
a_{\text{eff}}(s) = v^2 \kappa(s) - g \tan(\phi(s)).
\]
\end{definition}

\begin{constraint}
\[
|a_{\text{eff}}(s)| \leq a_{\text{eff,max}}.
\]
\end{constraint}

---

\subsection{Jerk}

\begin{definition}[Jerk]
The lateral jerk is defined as
\[
j(s) = \frac{\dd}{\dd t} a_y(s).
\]
\end{definition}

\begin{remark}
Using $s$ as parameter, one obtains
\[
j(s) = v^3 \kappa'(s).
\]
\end{remark}

\begin{constraint}
\[
|j(s)| \leq j_{\max}.
\]
\end{constraint}

---

\section{Coupling of Curvature and Cant}

\begin{remark}
Curvature and cant must not be treated independently.
\end{remark}

\begin{constraint}
\[
|\kappa'(s)| \leq f_1(v,\phi),
\qquad
|\phi'(s)| \leq f_2(v,\kappa).
\]
\end{constraint}

\begin{remark}
These relations express that:
\begin{itemize}
\item curvature changes require corresponding cant development,
\item cant changes must respect dynamic limits.
\end{itemize}
\end{remark}

---

\section{Speed-Dependent Constraints}

\begin{remark}
All relevant constraints depend on speed.
\end{remark}

\begin{definition}[Speed profile]
A speed profile is a function
\[
v : [0,L] \to \R_+.
\]
\end{definition}

\begin{remark}
Constraints must be evaluated pointwise with respect to $v(s)$.
\end{remark}

---

\section{Constraint Representation in AIM}

\subsection{Abstract Form}

\begin{definition}[Constraint functional]
A constraint is represented as a functional
\[
C_i[\kappa,\phi,v] \leq 0.
\]
\end{definition}

\begin{remark}
This allows constraints to be:
\begin{itemize}
\item evaluated continuously,
\item discretized for numerical optimization,
\item combined into a unified constraint system.
\end{itemize}
\end{remark}

---

\subsection{Discrete Representation}

\begin{remark}
In numerical implementation, constraints are evaluated on a grid
\[
s_0, \dots, s_N.
\]
\end{remark}

\begin{definition}[Discrete constraint]
\[
C_i(\mathbf{x}) = \max_k C_i(s_k).
\]
\end{definition}

---

\section{Integration into Optimization}

\begin{remark}
Constraints enter the optimization problem as:
\[
g_i(\mathbf{x}) \leq 0.
\]
\end{remark}

\begin{remark}
They may be handled via:
\begin{itemize}
\item hard constraints (feasibility),
\item penalty terms,
\item barrier methods.
\end{itemize}
\end{remark}

---

\section{Relation to Standards}

\begin{remark}
The presented constraints correspond to rules found in:
\begin{itemize}
\item UIC codes,
\item national railway regulations,
\item operator guidelines.
\end{itemize}
\end{remark}

\begin{remark}
The AIM framework provides a mathematical representation that allows
these rules to be applied consistently across different data sources.
\end{remark}

---

\section{Discussion}

\begin{remark}
Traditional workflows treat rules as external checks applied after
design.
\end{remark}

\begin{remark}
In contrast, AIM integrates rules directly into:
\begin{itemize}
\item modelling,
\item reconstruction,
\item optimization.
\end{itemize}
\end{remark}

\begin{remark}
This leads to a shift from validation to constraint-driven design.
\end{remark}

---

\section{Summary}

\begin{summary}
Railway design rules define the feasible set of alignment solutions.

Within the AIM framework, these rules are formalized as mathematical
constraints and integrated into the optimization process.

This allows alignment design to be treated as a constrained optimization
problem that respects both geometric and dynamic requirements.
\end{summary}
